\section{Постановка задачи}

В рамках данной работы необходимо разработать приложение телефонного справочника с использованием фреймворка Qt. Программа должна обеспечивать удобный графический интерфейс для работы с контактами и защищать пользователя от некорректного ввода данных.

Цель данной работы~--- реализовать программу, которая:
\begin{itemize}
    \item хранит контактную информацию: имя, фамилию, отчество, адрес, дату рождения, email и телефонные номера различных типов;
    \item выполняет валидацию всех вводимых данных с использованием регулярных выражений;
    \item позволяет добавлять, редактировать и удалять записи;
    \item осуществляет поиск по нескольким полям и сортировку по выбранному столбцу;
    \item сохраняет данные в файл и загружает их при запуске.
\end{itemize}

\subsection{Требования к валидации данных}

\subsubsection*{Фамилия, имя, отчество}
Поля ФИО должны удовлетворять следующим требованиям:
\begin{itemize}
    \item содержать только буквы различных алфавитов, цифры, дефис и пробел;
    \item начинаться с заглавной буквы;
    \item не начинаться и не заканчиваться дефисом;
    \item незначащие пробелы в начале и конце строки должны удаляться.
\end{itemize}

\subsubsection*{Телефонный номер}
Телефон должен быть записан в международном формате и храниться как последовательность цифр. Допустимые форматы ввода:
\begin{itemize}
    \item \texttt{+7 812 1234567}
    \item \texttt{8(800)123-1212}
    \item \texttt{+38 (032) 123 11 11}
\end{itemize}

Дополнительные ограничения: номер должен содержать от 7 до 15 цифр.

\subsubsection*{Дата рождения}
Дата рождения должна удовлетворять условиям:
\begin{itemize}
    \item быть меньше текущей даты;
    \item месяц~--- от 1 до 12;
    \item день~--- от 1 до 31 с учётом количества дней в месяце;
    \item учитывать високосные годы для февраля.
\end{itemize}

\subsubsection*{Адрес электронной почты}
Email должен соответствовать формату:
\begin{itemize}
    \item имя пользователя из латинских букв, цифр и спецсимволов;
    \item символ~\texttt{@};
    \item доменное имя из латинских букв, цифр и точек.
\end{itemize}

Все пробелы, включая пробелы вокруг символа \texttt{@}, должны быть удалены.

\subsection{Требования к функциональности}

\subsubsection*{Управление записями}
Приложение должно обеспечивать:
\begin{itemize}
    \item добавление новых контактов через диалоговое окно;
    \item редактирование существующих записей;
    \item удаление выбранных контактов с подтверждением.
\end{itemize}

\subsubsection*{Поиск и сортировка}
Необходимо реализовать:
\begin{itemize}
    \item поиск по всем полям или по конкретному полю;
    \item сортировку таблицы по любому столбцу;
    \item циклическое переключение направления сортировки: по возрастанию, по убыванию, без сортировки.
\end{itemize}

\subsubsection*{Хранение данных}
Данные должны сохраняться в файл при закрытии приложения и загружаться при запуске. Дополнительно требуется реализовать экспорт и импорт данных в формате \verb|YAML|.

\subsection{Требования к интерфейсу}

Для отображения данных используется табличное представление с колонками: фамилия, имя, телефон, email, адрес, дата рождения. Интерфейс должен включать:
\begin{itemize}
    \item главное окно с таблицей контактов;
    \item панель поиска с выбором поля для фильтрации;
    \item кнопки управления записями;
    \item меню с дополнительными функциями.
\end{itemize}